% ECONOMICS_PROBLEM: {"id": "G28-283", "title": "Net macroprudential benefits", "jel_code": "G28", "source_id": "OP-0324"}
% !TeX program = xelatex
\documentclass[11pt,a4paper]{article}
\usepackage[margin=23mm]{geometry}
\usepackage{fontspec}
\setmainfont{Latin Modern Roman}
\usepackage{amsmath,amssymb,mathtools}
\usepackage{xcolor,enumitem,booktabs,longtable,array,xurl,hyperref}
\definecolor{muted}{HTML}{53636C}
\hypersetup{colorlinks=true,urlcolor=blue,linkcolor=blue}
\setlength{\parindent}{0pt}
\setlength{\parskip}{5pt}
\newcommand{\R}{\mathbb R}
\newcommand{\E}{\mathbb E}
\newcommand{\N}{\mathbb N}
\newcommand{\ind}{\mathbf 1}
\newcommand{\Var}{\operatorname{Var}}
\newcommand{\Cov}{\operatorname{Cov}}
\newcommand{\supp}{\operatorname{supp}}
\DeclareMathOperator*{\argmin}{arg\,min}
\DeclareMathOperator*{\argmax}{arg\,max}
\newcommand{\entrymeta}[3]{{\sffamily\small\color{muted}\textbf{\texttt{#1}}\quad|\quad #2\par #3\par}}
\newcommand{\Problemref}[1]{\texttt{#1}}
\newcommand{\JELref}[2]{\texttt{#2} (JEL \texttt{#1})}
\begin{document}
\section*{Net macroprudential benefits}
\textbf{Problem ID:} \texttt{G28-283}\quad \textbf{JEL:} \texttt{G28}\par
% BEGIN ECONOMICS STATEMENT
\label{problem:OP-0324}
{\sffamily\small\color{muted}Original subject group: Banking, credit and financial stability\par}
\label{definitions:problem:30007606}
\textbf{Audit classification:} Published research agenda. The IMF source supports costs/resilience research. The household lifetime-welfare functional is editorial and is not identified by credit-growth estimates alone.
\subsection*{Definitions and precise research target}
Fix a target population of households with welfare weights \(\omega_i\geq0\), \(\sum_i\omega_i=1\), a discount factor \(\beta\in(0,1)\), and a feasible macroprudential policy \(a\). Household consumption under intervention \(do(a)\) is \(c_{it}(a)\); crisis losses, lending restrictions and fiscal costs enter the specified equilibrium rather than being added twice. Define
\[W(a)=E\left[\sum_{t=0}^{\infty}\beta^t\sum_i\omega_i u_i(c_{it}(a))\right],\]
assuming convergence and independently specified utility functions \(u_i\). The task is to identify or bound \(W(a)-W(0)\), its distribution across household groups, and the policy-induced change in crisis probability. Permit policy effects to depend on initial credit conditions and interact with monetary and microprudential rules. State how outcomes absent adoption are identified when regulators respond to latent risk. Account for nonbank and foreign lending, housing costs and borrowing access across income groups. A reduction in credit growth is neither a welfare benefit nor a welfare cost without these channels. Report sensitivity to welfare weights, tail-loss assumptions and discounting, and distinguish short-run causal estimates from long-run structural extrapolation.

\textbf{Additional formal definitions and scope (editorial).} Take a fixed finite household population or a probability measure over types, with integrable utility and \(\sum_t\beta^t E|\sum_i\omega_i u_i(c_{it}(a))|<\infty\). Include labor disutility and housing services if they change; otherwise state the restriction excluding them. Utility cardinalizations and welfare weights are fixed before cross-person aggregation. A crisis is a measurable event \(C(a)\) defined by a prespecified loss or financial-service interruption threshold over a fixed horizon. Use \(\Pr(C(a))-\Pr(C(0))\) separately from \(W(a)-W(0)\). A utility change need not be positive whenever crisis probability falls. For money-metric distributional reporting, define each type's compensating transfer through its baseline indirect utility, with monotonicity and range assumptions ensuring the transfer exists.

For the identification part, declare an admissible class \(\Theta\) of structural models, including preferences, technologies, shock laws, measurement/sampling rules and any equilibrium selector. If \(O\) denotes the complete observable history and \(T(\theta)\) the target defined above, the identified set is \(\mathcal I_T=\{T(\theta):\theta\in\Theta,\ \mathcal L_\theta(O)=\mathcal L_{\rm obs}(O)\}\); point identification means this set is a singleton. A \(do(a)\) comparison replaces stated policy/technology equations while fixing the declared exogenous law and re-solving the remaining equations. These definitions do not assert that an identification theorem holds without further restrictions.
\subsection*{Variants and assumption changes}
\begin{enumerate}
\item \textbf{Short-horizon policy evaluation (source-motivated).} Identify consumption/output and crisis-risk responses at fixed horizons under a credible policy assignment design.
\item \textbf{Lifetime distributional welfare (editorial).} Embed the responses in heterogeneous-household preferences and financing, then bound welfare under uncertain tail risks and weights.
\end{enumerate}
\subsection*{References and source audit}
\begin{enumerate}
\item Primary IMF DP 2023/002 cover confirms all seven authors; official landing-page author field omits Gelos and Martinez Peria. DOI verified at official metadata. Locator: Executive summary, printed pp.v--vi; Introduction, pp.viii--ix; Chapters 3--4 and 6. Support: Macroprudential costs, resilience and leakage agenda. PDF accessible at legacy IMF ashx URL; some repeated fetches throttled. URL: \url{https://www.imf.org/-/media/Files/Publications/DP/2023/English/MPEEOQEA.ashx}.
\end{enumerate}
\begin{enumerate}\raggedright
\item\hypertarget{definitions:source:30007606:1}{} Nina Biljanovska; Sophia Chen; Gaston Gelos; Deniz Igan; Maria Soledad Martinez Peria; Erlend Nier; Fabián Valencia. 2023. \emph{Macroprudential Policy Effects: Evidence and Open Questions}. Executive summary, printed pp.v--vi; Introduction, pp.viii--ix; Chapters 3--4 and 6.
\url{https://www.imf.org/-/media/Files/Publications/DP/2023/English/MPEEOQEA.ashx}.
DOI: \url{https://doi.org/10.5089/9798400226304.087}.
\end{enumerate}

\subsection*{Common conventions: Source mathematical conventions}

These conventions apply to each entry unless explicitly replaced.

\textbf{Models and identification.} A primitive model \(m\) specifies all probability laws, feasible choices, information, equilibrium and measurement rules used by its target. The admissible class \(\mathcal M\) is fixed independently of outcomes. If \(P_O(m)\) is its observable probability law and \(\tau(m)\) the target, its identified set at observable law \(P\) is
\[
 \mathcal I_\tau(P)=\{\tau(m):m\in\mathcal M,\ P_O(m)=P\}.
\]
Point identification means that this set is a singleton. An empty set signals incompatibility of data and assumptions. Coordinate bounds need not describe the joint set. Bounds are sharp only if every retained target is attainable or arbitrarily approachable by compatible models. Finite bounds require explicit restrictions on unobserved quantities; unrestricted confounding, hidden wealth, extrapolation or arbitrary equilibrium selection can yield uninformative sets.

\textbf{Causal interventions.} A potential outcome \(Y_i(p)\) is the outcome under a complete policy regime \(p\), including timing, financing, information and market spillovers. The target population and units are fixed before treatment. With interference, use the complete assignment vector or a justified exposure mapping; individual no-interference notation alone cannot identify an economy-wide intervention. A difference of mean potential outcomes does not require the cross-world joint distribution. Conditioning on post-treatment migration, survival, enrollment or employment changes the estimand and needs separate justification.

\textbf{Statistical statements.} An estimator, confidence set, forecast or test specifies a sampling law, model class and loss/coverage criterion. Uniform \((1-\alpha)\)-coverage means \(\inf_{m\in\mathcal M}\Pr_m\{\tau(m)\in C_n\}\geq1-\alpha\), or an explicitly asymptotic version. The whole parameter space trivially has coverage, so informativeness requires a stated width, risk or power objective. A forecasting improvement is relative to a named benchmark, information set and evaluation population.

\textbf{Equilibrium and optimization.} An equilibrium comprises feasible plans, optimality, beliefs where needed, budget constraints and all market/resource clearing. The primitives or a stated benchmark define each correspondence \(\mathcal E(p;m)\). Nonemptiness is assumed or proved, never inferred just from compact policy sets. A selected-equilibrium target uses a declared selection rule; a robust target quantifies over all permitted equilibria. Use suprema or infima unless attainment is established. An \(\varepsilon\)-optimal policy is feasible and within \(\varepsilon\) of the finite optimum value. Report an empty feasible set separately.

\textbf{Welfare and accounting.} Normative utility, interpersonal normalization, social weights, numeraire, discounting and the use of public receipts are primitives. Behavioral utility need not coincide with the welfare criterion. Household welfare includes ownership income and rents once; adding profits again to compensating variation generally double-counts them. A monetary equivalent is defined by an explicit transfer and price convention, with monotonicity and range conditions. Average effects, mechanism contributions, transfers and welfare losses are different targets. Infinite sums require convergence and dynamic feasible plans require terminal or no-Ponzi conditions.

\textbf{Source versions.} A working-paper date, revision date and journal publication year are distinguished. A DOI for a journal version is not automatically the DOI of an earlier manuscript. Locators refer to the stated version and printed pages where specified. A metadata-only check does not verify a claim about a full-text conclusion. Every entry's source subsection narrows the provenance claim accordingly.
% END ECONOMICS STATEMENT
\end{document}
